\documentclass[11pt,twoside,leqno]{article}

% =============================
% Packages (base user preference)
% =============================
\usepackage[T1]{fontenc}
\usepackage[utf8]{inputenc}
\usepackage{lmodern}
\usepackage[a4paper,margin=1in]{geometry}

\usepackage{amsmath,amsfonts,amssymb,amsthm,mathtools}
\usepackage{enumitem}
\usepackage{booktabs,array}
\usepackage{graphicx}

\usepackage{tikz}
\usetikzlibrary{arrows.meta,positioning}

\usepackage{stmaryrd}

\usepackage{algorithm}
\usepackage{algpseudocode}

\usepackage[colorlinks=true,linkcolor=blue,citecolor=blue,urlcolor=blue]{hyperref}

% =============================
% Theorem environments
% =============================
\theoremstyle{definition}
\newtheorem{definition}{Definition}[section]
\newtheorem{example}[definition]{Example}

\theoremstyle{plain}
\newtheorem{theorem}[definition]{Theorem}
\newtheorem{lemma}[definition]{Lemma}
\newtheorem{proposition}[definition]{Proposition}
\newtheorem{corollary}[definition]{Corollary}

\theoremstyle{remark}
\newtheorem{remark}[definition]{Remark}

% =============================
% User macros (kept as requested)
% =============================
\newcommand{\Atr}{\mathsf{Atr}}
\newcommand{\Pre}{\mathsf{Pre}}
\newcommand{\CPre}{\mathsf{CPre}}
\newcommand{\lfp}{\mathsf{lfp}}
\newcommand{\fun}{\mathsf{fun}}
\newcommand{\Pow}{\mathcal{P}}
\newcommand{\N}{\mathbb{N}}

% =============================
% Additional macros for this paper
% =============================
\newcommand{\F}{\mathbb{F}}
\newcommand{\Ftwo}{\mathbb{F}_2}
\newcommand{\wt}{\mathsf{wt}}
\newcommand{\supp}{\mathsf{supp}}
\newcommand{\syn}{\mathsf{syn}}
\newcommand{\rk}{\mathsf{rk}}
\newcommand{\HH}{\mathbf{H}}
\newcommand{\ee}{\mathbf{e}}
\newcommand{\xx}{\mathbf{x}}
\newcommand{\cc}{\mathbf{c}}
\newcommand{\diam}{\mathsf{diam}}
\newcommand{\covrad}{\rho}

\title{A Canonical Poset--TDA Pipeline for Binary Linear Codes\\
via the Syndrome Cayley Graph and Rank--Truncated Order Complexes}
\author{Carlos Ernesto Ram\'irez Ovalle}
\date{\today}

\begin{document}
\maketitle

\begin{abstract}
We present a canonical topological data analysis (TDA) pipeline for binary linear codes that replaces
coset-leader-dependent constructions by an intrinsic syndrome-level framework. Given a parity-check matrix
$\HH$, we consider the Cayley graph of the syndrome space generated by the columns of $\HH$ and define the
\emph{syndrome rank} $\rk(s)$ as the graph distance from $0$. This rank coincides with the minimum Hamming
weight among all error patterns realizing $s$ (syndrome weight). From adjacency and rank increments we define a
canonical cover relation and obtain a graded poset on syndromes. We then build a filtration of simplicial complexes
from rank-truncated subposets via order complexes and compute persistent homology. We provide self-contained proofs
of well-definedness, gradedness, and invariance under row operations and coordinate permutations, and we explain
why the unpunctured order-complex filtration is necessarily contractible (hence yields $\beta_1=0$), motivating
punctured/relative variants as the informative TDA objects. We also outline an exact computational regime and a
reproducible algorithmic workflow suitable for small to moderate redundancy $m=n-k$.
\end{abstract}

\paragraph{Keywords.}
Binary linear codes; coset leaders; syndrome decoding; Cayley graphs; poset topology; order complexes; persistent homology.



% ==========================================================
\section{Introduction}\label{sec:introduction}
% ==========================================================

\subsection{Motivation: from coset-leader Hasse diagrams to canonical syndrome structures}
A classical way to organize the cosets of a binary linear code is to select a \emph{coset leader} for each coset
and draw a Hasse-type diagram whose cover relations reflect adding (or removing) one coordinate in the support of
a minimum-weight representative. This perspective is useful for decoding and visualization, but it depends on a
choice of leaders, and leaders are not unique in general. Consequently, two valid leader-selection strategies may
produce different adjacency information even for the same code.

The main idea of this paper is to \emph{canonize} the construction at the syndrome level: instead of picking leaders,
we use the full set of columns of a parity-check matrix to define a Cayley graph on the syndrome space and derive
a canonical rank, cover relation, and filtration for TDA.

\subsection{Contributions and paper structure (the dotted list)}
The following dotted list summarizes the paper's contributions and also serves as a reading map. Each item is a
self-contained block of results; the order matches the logical dependencies.

\begin{itemize}[leftmargin=*,itemsep=2pt]
\item[\textbullet] \textbf{Canonical syndrome rank (Section~\ref{sec:rank}).}
We define $\rk(s)$ as the distance from $0$ to $s$ in the syndrome Cayley graph generated by the columns of $\HH$
and prove $\rk(s)$ equals the minimum Hamming weight among all error patterns with syndrome $s$.

\item[\textbullet] \textbf{Canonical graded poset on syndromes (Section~\ref{sec:poset}).}
We define the cover relation using adjacency plus rank increment and prove it induces a graded poset with rank function $\rk$.

\item[\textbullet] \textbf{Invariance (Section~\ref{sec:invariance}).}
We prove the construction is invariant under row operations on $\HH$ (change of parity-check representation) and coordinate permutations.

\item[\textbullet] \textbf{Filtration 2 via order complexes (Section~\ref{sec:filtration}).}
We build a simplicial filtration $K_t=\Delta(P_{\le t})$ from rank-truncated subposets and show why the unpunctured filtration is a cone
(and hence topologically trivial), motivating punctured/relative variants as the correct objects for informative persistence.

\item[\textbullet] \textbf{Algorithms and exact parameter regime (Sections~\ref{sec:algorithms}--\ref{sec:regime}).}
We give a BFS-based exact computation of $\rk$ with complexity $O(n2^m)$ and discuss scientifically reasonable truncations for
building simplicial complexes up to fixed dimension.

\item[\textbullet] \textbf{Sanity checks and interpretation of anomalies (Section~\ref{sec:sanity}).}
We connect $\max_s\rk(s)$ to the covering radius and explain how perfect-code constraints (e.g.\ Golay) provide validation tests.
\end{itemize}

\subsection{Related work (positioning)}\label{sec:related}
Graphical representations of binary linear codes via coset leaders and Hasse diagrams have been developed and studied in the coding-theory
literature; see, e.g., \cite{DelgadoCastilloHolguin2022} and the thesis-level development \cite{DelgadoThesis2021}. The notion that the minimal
coset weight corresponds to a minimal way to write a syndrome as a linear combination of columns of a parity-check matrix appears explicitly in
works on coset leader weight enumerators and syndrome weight \cite{Blokhuis2022,PellikaanSlides}. Poset topology and order complexes are standard
tools in algebraic and topological combinatorics \cite{BjornerTopMeth,WachsPosetTopology}. For persistent homology and computational TDA we use the
standard framework \cite{EdelsbrunnerHarer2010,ZomorodianCarlsson2005}, and for data structures we refer to the simplex tree \cite{BoissonnatMaria2014}.

% ==========================================================
\section{Preliminaries}\label{sec:prelims}
% ==========================================================

\subsection{Binary linear codes, cosets, and syndromes}
Let $C\le \Ftwo^n$ be an $[n,k]$ binary linear code. Fix a parity-check matrix
$\HH\in \Ftwo^{m\times n}$ with $m=n-k$ such that $C=\ker(\HH)$.
Write $\HH=[h_1|\cdots|h_n]$ with columns $h_i\in\Ftwo^m$.

\begin{definition}[Syndrome map and fibers]\label{def:syndrome}
Define $\syn:\Ftwo^n\to \Ftwo^m$ by $\syn(\ee)=\HH \ee^\top$.
For $s\in \Ftwo^m$, define the syndrome class (fiber)
\[
E_s\coloneqq\{\ee\in \Ftwo^n:\syn(\ee)=s\}.
\]
\end{definition}

\begin{lemma}[Syndromes parametrize cosets]\label{lem:cosets}
For each $s\in\Ftwo^m$ and any $\ee\in E_s$, we have $E_s=\ee+C$.
Moreover, $\{E_s\}_{s\in\Ftwo^m}$ partitions $\Ftwo^n$ into exactly $2^m$ cosets.
\end{lemma}

\begin{proof}
If $\cc\in C$, then $\syn(\ee+\cc)=\HH(\ee+\cc)^\top=\HH\ee^\top+\HH\cc^\top=s$,
so $\ee+C\subseteq E_s$. Conversely, if $\ee'\in E_s$, then
$\HH(\ee-\ee')^\top=0$, hence $\ee-\ee'\in C$ and $\ee'\in \ee+C$. Thus $E_s=\ee+C$.
Distinct syndromes yield disjoint fibers, and $|\Ftwo^m|=2^m$.
\end{proof}

\subsection{Support and Hamming weight}
\begin{definition}[Support and Hamming weight]\label{def:weight}
For $\ee=(e_1,\dots,e_n)\in\Ftwo^n$ define
$\supp(\ee)=\{i: e_i=1\}$ and $\wt(\ee)=|\supp(\ee)|$.
\end{definition}

\begin{definition}[Syndrome weight / coset-leader weight]\label{def:syndromewt}
For $s\in\Ftwo^m$, define the syndrome weight
\[
\wt_{\HH}(s)\coloneqq \min\{\wt(\ee): \ee\in E_s\}.
\]
\end{definition}

\begin{remark}
The minimum in Definition~\ref{def:syndromewt} need not be attained by a unique vector.
Any $\ee\in E_s$ with $\wt(\ee)=\wt_{\HH}(s)$ is a coset leader (not necessarily unique).
\end{remark}

% ==========================================================
\section{Syndrome Cayley graph and canonical rank}\label{sec:rank}
% ==========================================================

\begin{definition}[Syndrome Cayley graph]\label{def:cayley}
Define the undirected graph $G_\HH$ with vertex set $\Ftwo^m$ and edges
\[
s \sim s+h_i \qquad (s\in\Ftwo^m,\ i\in\{1,\dots,n\}),
\]
where addition is in $\Ftwo^m$. Let $d_{G_\HH}$ denote shortest-path distance in $G_\HH$.
Define the canonical rank function
\[
\rk(s)\coloneqq d_{G_\HH}(0,s).
\]
\end{definition}

\begin{lemma}[Syndrome as a sum of columns]\label{lem:syndrome-sum}
For $\ee\in\Ftwo^n$ with support $S=\supp(\ee)$,
\[
\syn(\ee)=\HH\ee^\top=\sum_{i\in S} h_i \in \Ftwo^m.
\]
\end{lemma}
\begin{proof}
Over $\Ftwo$, $\HH\ee^\top=\sum_{i=1}^n e_i h_i$ and $e_i\in\{0,1\}$, hence the sum reduces to $i\in S$.
\end{proof}

\begin{theorem}[Rank equals syndrome weight]\label{thm:rk-equals-wtH}
For every $s\in\Ftwo^m$,
\[
\rk(s)=\wt_{\HH}(s)=\min\{\wt(\ee): \ee\in\Ftwo^n,\ \HH\ee^\top=s\}.
\]
\end{theorem}

\begin{proof}
\noindent\textbf{(1) $\rk(s)\le \wt_{\HH}(s)$.}
Take $\ee\in E_s$ with $\wt(\ee)=\wt_{\HH}(s)$ and let $S=\supp(\ee)$, $|S|=\wt_{\HH}(s)$.
By Lemma~\ref{lem:syndrome-sum}, $s=\sum_{i\in S}h_i$.
Consider the walk in $G_\HH$:
\[
0\to h_{i_1}\to h_{i_1}+h_{i_2}\to\cdots\to \sum_{j=1}^{|S|}h_{i_j}=s,
\]
where $(i_1,\dots,i_{|S|})$ is any ordering of $S$. This is a path of length $|S|$, hence $\rk(s)\le |S|=\wt_{\HH}(s)$.

\smallskip
\noindent\textbf{(2) $\wt_{\HH}(s)\le \rk(s)$.}
Let $L=\rk(s)=d_{G_\HH}(0,s)$ and choose a shortest path
$0=s_0\sim s_1\sim\cdots\sim s_L=s$.
For each step there exists $i_j$ such that $s_j=s_{j-1}+h_{i_j}$.
Define $\ee\in\Ftwo^n$ by setting $e_i=1$ iff $i$ appears an odd number of times among $\{i_1,\dots,i_L\}$.
Then $\wt(\ee)\le L$, and
\[
\HH\ee^\top=\sum_{j=1}^L h_{i_j}=\sum_{j=1}^L(s_{j-1}+s_j)=s_0+s_L=0+s=s,
\]
so $\ee\in E_s$ and $\wt_{\HH}(s)\le \wt(\ee)\le L=\rk(s)$.

Combining (1) and (2) yields $\rk(s)=\wt_{\HH}(s)$.
\end{proof}

\begin{theorem}[Characterization of cosets via syndromes]\label{thm:coset-syndrome-en}
Let $C\le \mathbb{F}_2^n$ be a binary linear code and let
$H\in \mathbb{F}_2^{(n-k)\times n}$ be a parity-check matrix of $C$, that is,
\[
C=\{v\in\mathbb{F}_2^n:\; vH^{\mathsf T}=0\}.
\]
Define the \emph{syndrome} map $S:\mathbb{F}_2^n\to \mathbb{F}_2^{n-k}$ by
\[
S(x)=xH^{\mathsf T}.
\]
Then for any $x,y\in \mathbb{F}_2^n$ the following statements are equivalent:
\begin{enumerate}[label=\textnormal{(\roman*)}]
\item $S(x)=S(y)$.
\item $y-x\in C$ \;\;(equivalently, $y+x\in C$, since over $\mathbb{F}_2$ we have $x=-x$).
\item $x$ and $y$ belong to the same coset of $C$, i.e., $x+C=y+C$.
\end{enumerate}
In particular, the syndrome depends \emph{only} on the coset: if $x+C=y+C$ then $S(x)=S(y)$.
\end{theorem}

\begin{proof}
We first record two basic facts.

\smallskip
\noindent\textbf{(A) Linearity of $S$.}
For any $u,v\in\mathbb{F}_2^n$,
\[
S(u+v)=(u+v)H^{\mathsf T}=uH^{\mathsf T}+vH^{\mathsf T}=S(u)+S(v),
\]
since matrix multiplication distributes over addition in the field $\mathbb{F}_2$.
Moreover, for $\lambda\in\mathbb{F}_2$,
\[
S(\lambda u)=(\lambda u)H^{\mathsf T}=\lambda(uH^{\mathsf T})=\lambda S(u),
\]
hence $S$ is a linear map.

\smallskip
\noindent\textbf{(B) Kernel of $S$.}
Because $H$ is a parity-check matrix for $C$, we have
\[
x\in C\quad\Longleftrightarrow\quad xH^{\mathsf T}=0
\quad\Longleftrightarrow\quad S(x)=0.
\]
Thus $\ker(S)=C$.

\smallskip
We now prove \textnormal{(i)}$\Leftrightarrow$\textnormal{(ii)}$\Leftrightarrow$\textnormal{(iii)}.

\smallskip
\noindent\textbf{\textnormal{(i)} $\Rightarrow$ \textnormal{(ii)}.}
Assume $S(x)=S(y)$. Using linearity \textbf{(A)},
\[
S(y-x)=S(y)+S(-x).
\]
But over $\mathbb{F}_2$ we have $-x=x$, hence $S(-x)=S(x)$ and therefore
\[
S(y-x)=S(y)+S(x)=0,
\]
because $S(y)=S(x)$ implies $S(y)+S(x)=0$ in $\mathbb{F}_2^{n-k}$.
By \textbf{(B)}, $S(y-x)=0$ implies $y-x\in C$, which is \textnormal{(ii)}.

\smallskip
\noindent\textbf{\textnormal{(ii)} $\Rightarrow$ \textnormal{(iii)}.}
Assume $y-x\in C$. Then there exists $c\in C$ such that $y-x=c$, i.e., $y=x+c$.
Consider the coset of $y$:
\[
y+C=(x+c)+C=x+(c+C).
\]
Since $c\in C$ and $C$ is a subgroup under addition, we have $c+C=C$.
Hence
\[
y+C=x+C,
\]
which is exactly \textnormal{(iii)}.

\smallskip
\noindent\textbf{\textnormal{(iii)} $\Rightarrow$ \textnormal{(i)}.}
Assume $x+C=y+C$. By the definition of coset equality, this implies $y\in x+C$,
so there exists $c\in C$ with $y=x+c$.
Applying $S$ and using linearity \textbf{(A)},
\[
S(y)=S(x+c)=S(x)+S(c).
\]
But $c\in C=\ker(S)$ by \textbf{(B)}, hence $S(c)=0$ and therefore
\[
S(y)=S(x).
\]
This proves \textnormal{(i)}.

\smallskip
We conclude that \textnormal{(i)}, \textnormal{(ii)} and \textnormal{(iii)} are equivalent.
The final statement (“the syndrome depends only on the coset”) is precisely \textnormal{(iii)}$\Rightarrow$\textnormal{(i)}.
\end{proof}

\begin{remark}[Computational consequence]
By Theorem~\ref{thm:rk-equals-wtH}, $\rk(\cdot)$ is computed exactly by a BFS from $0$ in $G_\HH$ with time $O(n2^m)$ and memory $O(2^m)$.
\end{remark}

\subsection{Covering radius and the maximum rank}
\begin{definition}[Covering radius]\label{def:coveringradius}
The covering radius of a code $C\le \Ftwo^n$ is
\[
\covrad(C)\coloneqq \max_{\xx\in\Ftwo^n}\min_{\cc\in C}\wt(\xx-\cc).
\]
\end{definition}

\begin{proposition}[Maximum syndrome weight equals covering radius]\label{prop:rho}
For a binary linear code with parity-check matrix $\HH$,
\[
\covrad(C)=\max_{s\in\Ftwo^m}\wt_{\HH}(s)=\max_{s\in\Ftwo^m}\rk(s).
\]
\end{proposition}

\begin{proof}
Every $\xx\in\Ftwo^n$ lies in a unique coset $\xx+C=E_{\syn(\xx)}$ (Lemma~\ref{lem:cosets}).
Moreover, $\min_{\cc\in C}\wt(\xx-\cc)$ is the minimum weight in that coset, i.e.\ $\wt_{\HH}(\syn(\xx))$.
Taking the maximum over $\xx$ is therefore the same as taking the maximum over syndromes $s$.
Finally, $\wt_{\HH}(s)=\rk(s)$ by Theorem~\ref{thm:rk-equals-wtH}.
\end{proof}

% ==========================================================
\section{Canonical cover relation and graded poset}\label{sec:poset}
% ==========================================================

\begin{definition}[Cover relation]\label{def:cover}
For $u,v\in\Ftwo^m$, define $u\prec v$ if there exists $i\in\{1,\dots,n\}$ such that
\[
v=u+h_i
\quad\text{and}\quad
\rk(v)=\rk(u)+1.
\]
\end{definition}

\begin{definition}[Induced partial order]\label{def:order}
Define $u\le v$ if either $u=v$, or there exists a finite chain
$u=u_0\prec u_1\prec\cdots\prec u_\ell=v$.
\end{definition}

\begin{theorem}[Gradedness]\label{thm:graded}
$(\Ftwo^m,\le)$ is a graded poset with rank function $\rk$. In particular:
(i) if $u\prec v$ then $\rk(v)=\rk(u)+1$; and
(ii) if $u<v$ then $\rk(u)<\rk(v)$.
\end{theorem}

\begin{proof}
Reflexivity and transitivity are immediate from Definition~\ref{def:order}.
For antisymmetry, assume $u\le v$ and $v\le u$ with $u\neq v$.
Then there exists a nontrivial chain $u=u_0\prec\cdots\prec u_\ell=v$, $\ell\ge 1$.
By Definition~\ref{def:cover}, $\rk(u_j)=\rk(u)+j$, so $\rk(v)=\rk(u)+\ell>\rk(u)$, contradicting $v\le u$.
Hence $u=v$.

Gradedness follows because covers increase rank by $1$, and rank strictly increases along strict chains.
\end{proof}

% ==========================================================
\section{Invariance properties}\label{sec:invariance}
% ==========================================================

\subsection{Row-operation invariance}
\begin{theorem}[Row-operation invariance]\label{thm:rowinv}
Let $\HH'\in\Ftwo^{m\times n}$ satisfy $\HH'=U\HH$ for some invertible $U\in GL(m,\Ftwo)$.
Then $G_{\HH}$ and $G_{\HH'}$ are isomorphic via $\varphi:\Ftwo^m\to\Ftwo^m$, $\varphi(s)=Us$.
Consequently, $\rk_{\HH'}(\varphi(s))=\rk_{\HH}(s)$ for all $s$, and the induced posets are isomorphic.
\end{theorem}

\begin{proof}
Columns of $\HH'$ are $h'_i=Uh_i$. An edge in $G_\HH$ has the form $s\sim s+h_i$.
Applying $\varphi$ gives $\varphi(s)\sim \varphi(s)+Uh_i=\varphi(s)+h'_i$, an edge of $G_{\HH'}$.
Since $U$ is invertible, $\varphi$ is a bijective graph isomorphism, hence preserves distances.
The cover relation is adjacency plus rank increment, thus preserved by $\varphi$.
\end{proof}

\subsection{Coordinate-permutation invariance}
\begin{theorem}[Column-permutation invariance]\label{thm:colperm}
If $\HH'$ is obtained from $\HH$ by permuting columns, then $G_{\HH'}=G_{\HH}$ as unlabeled graphs.
Hence $\rk$ and the poset structure are identical.
\end{theorem}

\begin{proof}
Permuting columns permutes the multiset of generators $\{h_i\}$, leaving the adjacency rule $s\sim s+h_i$ unchanged.
Therefore the graph and all distances are unchanged.
\end{proof}

% ==========================================================
\section{Filtration 2: order complexes of rank truncations}\label{sec:filtration}
% ==========================================================

\subsection{Rank truncations and order complexes}
\begin{definition}[Rank truncation]
For $t\in\N$, define the induced subposet
\[
P_{\le t}\coloneqq\{s\in\Ftwo^m:\rk(s)\le t\}
\]
with order inherited from $(\Ftwo^m,\le)$.
\end{definition}

\begin{definition}[Order complex filtration]\label{def:ordercomplex}
Define
\[
K_t\coloneqq \Delta(P_{\le t}),
\]
where $\Delta(\cdot)$ is the order complex: simplices are chains $s_0<\cdots<s_k$ in $P_{\le t}$.
Then $K_0\subseteq K_1\subseteq\cdots$ is a simplicial filtration.
\end{definition}

\begin{lemma}[Membership by rank]\label{lem:filt}
A chain $\sigma=\{s_0<\cdots<s_k\}$ belongs to $K_t$ if and only if $\max_j\rk(s_j)\le t$.
\end{lemma}
\begin{proof}
$\sigma\in K_t$ iff all vertices lie in $P_{\le t}$ iff $\rk(s_j)\le t$ for all $j$.
\end{proof}

\subsection{Why the unpunctured filtration is topologically trivial}
The poset has a unique minimum element, namely $0\in\Ftwo^m$, and $0\in P_{\le t}$ for all $t$.
This forces each $K_t$ to be a cone, hence contractible.

\begin{proposition}[Cone property]\label{prop:cone}
For every $t$, the complex $K_t=\Delta(P_{\le t})$ is contractible.
Consequently, $\widetilde H_p(K_t)=0$ for all $p\ge 0$, and in particular $\beta_1(K_t)=0$.
\end{proposition}

\begin{proof}
Let $\sigma$ be any simplex of $K_t$, i.e.\ a chain in $P_{\le t}$. Since $0$ is the minimum element, $0<s$ for every $s\in P_{\le t}\setminus\{0\}$,
hence $\sigma\cup\{0\}$ is also a chain and therefore a simplex in $K_t$. Thus every simplex is contained in a simplex that includes $0$,
so $K_t$ is a cone with apex $0$ and is contractible.
\end{proof}

\subsection{Informative variants: punctured and relative filtrations}
Proposition~\ref{prop:cone} implies that the \emph{unpunctured} filtration cannot produce nontrivial $H_1$ (and more generally, no reduced homology).
Therefore, the informative TDA objects are obtained by removing the cone apex or by using relative homology.

\begin{definition}[Punctured filtration]\label{def:punctured}
Define the punctured complexes
\[
K_t^\circ \coloneqq \Delta\big(P_{\le t}\setminus\{0\}\big).
\]
\end{definition}

\begin{remark}[What changes after puncturing]
The complexes $K_t^\circ$ need not be connected; thus $\beta_0(K_t^\circ)$ and higher Betti numbers may vary with the code and with $t$.
In particular, observations of $\beta_0>1$ are meaningful only if the apex has been removed (or if the construction is incomplete due to truncation artifacts).
\end{remark}

% ==========================================================
\section{Algorithms and reproducibility}\label{sec:algorithms}
% ==========================================================

\subsection{Exact ranks by BFS}
\begin{algorithm}[h]
\caption{Canonical ranks by BFS in the syndrome Cayley graph}\label{alg:bfs}
\begin{algorithmic}[1]
\Require Columns $h_1,\dots,h_n\in\Ftwo^m$ of $\HH$
\Ensure $\rk(s)$ for all $s\in\Ftwo^m$
\State Initialize $\rk(0)\gets 0$ and $\rk(s)\gets -1$ for $s\neq 0$
\State Queue $Q\gets [0]$
\While{$Q$ not empty}
  \State pop $u$ from $Q$
  \For{$i=1$ to $n$}
     \State $v\gets u+h_i$
     \If{$\rk(v)=-1$}
        \State $\rk(v)\gets \rk(u)+1$
        \State push $v$ into $Q$
     \EndIf
  \EndFor
\EndWhile
\end{algorithmic}
\end{algorithm}

\subsection{Cover edges and simplices}
Given $\rk$, define the cover edges $u\prec v$ using Definition~\ref{def:cover}. To build $K_t$ (or $K_t^\circ$),
one can either:
\begin{enumerate}[leftmargin=*]
\item construct the Hasse diagram (cover graph) and enumerate chains up to a fixed dimension $d_{\max}$, inserting all faces; or
\item enumerate saturated chains (rank-increasing by $1$ at each step) up to a rank cutoff $t_{\max}$ and insert all faces up to dimension $d_{\max}$.
\end{enumerate}

\subsection{Persistent homology computation}
Once the filtered simplicial complex is stored (e.g.\ as a simplex tree), compute persistence diagrams/barcodes and Betti curves using standard TDA libraries.
We recommend reporting both (i) barcodes/diagrams and (ii) scalar summaries such as maximum persistence, number of long bars, and Betti curves as functions of $t$.

% ==========================================================
\section{Scientifically reasonable exact regime (binary)}\label{sec:regime}
% ==========================================================

The syndrome space has size $2^m=2^{n-k}$. Exact computation of $\rk$ is feasible for moderate $m$.

\begin{proposition}[Practical exact regime]\label{prop:regime}
A practical exact regime is $m=n-k\le 12$ (at most $4096$ syndromes), with simplex-dimension truncation
$d_{\max}\in\{2,3\}$ and rank cutoff $t_{\max}\le 6$ if needed to control simplex growth.
\end{proposition}

\begin{proof}
BFS requires time $O(n2^m)$ and memory $O(2^m)$. The number of simplices depends on the number of chains and can grow quickly;
fixing $d_{\max}$ and optionally truncating by $t_{\max}$ are standard and scientifically reasonable controls in computational topology.
\end{proof}

% ==========================================================
\section{Sanity checks and interpretation of anomalies}\label{sec:sanity}
% ==========================================================

\subsection{Perfect-code check (Golay as validation)}
If $C$ is a perfect $t$-error-correcting code, then every syndrome is realized by an error of weight $\le t$.
Equivalently, $\max_s \rk(s)=t$ (Proposition~\ref{prop:rho}). For the binary Golay perfect code $[23,12,7]$ we have $t=3$ and
\[
\sum_{i=0}^{3}\binom{23}{i}=2048=2^{11},
\]
so weight-$\le 3$ errors exhaust all $2^{11}$ syndromes. Therefore, a computation yielding $\max_s\rk(s)>3$ indicates either
(i) the parity-check matrix used is not equivalent to the Golay code, or (ii) there is a mismatch/bug in the generator set or data parsing.

\subsection{Why $\beta_1=0$ in Filtration 2 (unpunctured)}
By Proposition~\ref{prop:cone}, the unpunctured filtration $K_t=\Delta(P_{\le t})$ is always a cone and therefore has $\beta_1(K_t)=0$ for all $t$.
Hence empirical observations ``all codes have $\beta_1=0$'' are expected and should be reported as a \emph{theoretical sanity check},
not as a code-discriminating feature.

\subsection{When $\beta_0>1$ can appear}
In a correct unpunctured implementation, $K_t$ is connected and contractible, so $\beta_0(K_t)=1$ for all $t$.
Thus $\beta_0>1$ may occur only if:
\begin{enumerate}[leftmargin=*]
\item you are computing the punctured filtration $K_t^\circ$ (where multiple components are meaningful); or
\item the complex is incomplete due to truncation artifacts (e.g.\ enumerating chains by \emph{length} rather than by rank depth, or missing faces).
\end{enumerate}

% ==========================================================
\section{Conclusion}\label{sec:conclusion}
% ==========================================================

We provided a canonical syndrome-level poset--TDA framework for binary linear codes:
ranks via BFS on the syndrome Cayley graph, a canonical cover relation and graded poset, and Filtration~2 via
order complexes of rank truncations. The framework is invariant under row operations and coordinate permutations,
and it clarifies which homological signals are theoretically meaningful (punctured/relative variants) versus
structurally forced (cone triviality in the unpunctured filtration). This canonized viewpoint is designed to
improve and stabilize coset-leader-based graphical representations by removing non-uniqueness and enabling
reproducible persistence-based descriptors.

% ==========================================================
\begin{thebibliography}{99}

\bibitem{DelgadoThesis2021}
L.~D.~Delgado Ordo\~nez.
\newblock \emph{Proceso de decodificaci\'on en c\'odigos binarios usando l\'ideres de clases laterales}.
\newblock Undergraduate/degree thesis, Universidad del Cauca, Popay\'an, 2021.

\bibitem{DelgadoCastilloHolguin2022}
L.~Delgado, J.~H.~Castillo, and A.~Holgu\'in-Villa.
\newblock \emph{A graphical representation of binary linear codes}.
\newblock arXiv:2205.10574, 2022.

\bibitem{Blokhuis2022}
A.~Blokhuis, A.~E.~Brouwer, and others.
\newblock The extended coset leader weight enumerator of a twisted cubic code.
\newblock \emph{Designs, Codes and Cryptography}, 2022.

\bibitem{PellikaanSlides}
R.~Pellikaan.
\newblock \emph{The coset leader weight enumerator of the code of the twisted cubic} (slides).
\newblock CIRM, available online.

\bibitem{BjornerTopMeth}
A.~Bj{\"o}rner.
\newblock Topological methods.
\newblock In \emph{Handbook of Combinatorics}, 1995.

\bibitem{WachsPosetTopology}
M.~Wachs.
\newblock \emph{Poset Topology: Tools and Applications}.
\newblock IAS/Park City Mathematics Series lecture notes, 2004.

\bibitem{EdelsbrunnerHarer2010}
H.~Edelsbrunner and J.~Harer.
\newblock \emph{Computational Topology: An Introduction}.
\newblock AMS, 2010.

\bibitem{ZomorodianCarlsson2005}
A.~Zomorodian and G.~Carlsson.
\newblock Computing persistent homology.
\newblock \emph{Discrete \& Computational Geometry}, 33(2):249--274, 2005.

\bibitem{BoissonnatMaria2014}
J.-D.~Boissonnat and C.~Maria.
\newblock The simplex tree: an efficient data structure for general simplicial complexes.
\newblock \emph{Algorithmica}, 70:406--427, 2014.

\end{thebibliography}

\end{document}
